\section{电子-声子耦合}
在很多体系中，因为电子的运动比原子核要快得多，所以电子和晶格振动自由度(声子)可以分别处理。不过这是一种近似处理，必要时可以通过引入电子-声子耦合来修正计算结果。这就涉及到两个体系的耦合计算，尽管很大程度上，两个体系还是独立处理的。实际上，\textcolor{red}{电子-声子散射是广泛应用的普遍效应。}如室温下，半导体的迁移率或导体的电导率问题等。

\subsection{微扰理论的带隙重整化}
\subsubsection{金刚石的温度相关带隙重整化}
对于16-原子的金刚石超晶胞计算零点和温度相关的声子诱导的带隙重整化。

电子-声子耦合计算至少包含两个核心步骤。第一步计算中，通过超晶胞计算电子-声子相互作用势，第二步，通过初基原胞，计算由电子-声子相互作用引起的能带重整化。通过该计算流程，可以将复杂的高开销的超晶胞计算与最后元胞性质计算分离开来，达到精确、有效地计算电子-声子效应的目的。

作为本算例的第一步，首先通过超晶胞和有限差分方法，计算金刚石的电子-声子相互作用势。金刚石作为共价绝缘体存在宽的间接带隙，是研究电子-声子耦合非常好的模型体系。其内部强的共价键将会导致高频声子和显著的电子-声子耦合。理解其带隙重整化是精确预测不同温度下的光学和电学性质的重要前提。

本算例中相关的定量计算都以\textrm{HDF5}文件格式保存在文件\textsc{phelel\_params.hdf5}中。在本算例的第二步，采用上述电子-声子耦合势计算元胞中的$g_{mnv}(\vec k,\vec q)$，再来计算电子-声子耦合和由此引起的带隙重整化。最后，在零温和有限温度下计算金刚石的直接带隙和间接带隙(基础带隙)的重整化。

$2\times2\times2$金刚石结构\textrm{POSCAR}

\textrm{POSCAR}定义的是金刚石超晶胞结构。在本算例中，共有$2\times2\times2$的初基原胞的复本。之所以需要超晶胞，是因为计算中需要使用有限原子偏移。

同样出于计算开销的原因，这里选择相对较小的$2\times2\times2$超晶胞。在实际计算过程中，超晶胞的大小要审慎地选取，根据所计算的物理量和超晶胞大小的收敛测试结果来确定。如针对带隙重整化，依次采用$2\times2\times2$，$3\times3\times3$，$4\times4\times4$逐渐增大的超晶胞来做收敛测试。

计算控制参数文件\textrm{INCAR}

在\textrm{INCAR}文件中，有两个关于电子-声子相互作用势的重要计算控制参数
\begin{itemize}
	\item \texttt{IBRION}:~确定\textrm{VASP}计算中选取的具体的有限差分方法
	\item \texttt{ELPH\_POT\_GENERATE}:~指定计算电子-声子相互作用势
\end{itemize}

$\vec k$空间网格剖分文件\textrm{KPOINTS}，这里只选择$\Gamma$点。

进入计算目录，执行\textrm{Bash}命令，完成\textrm{VASP}计算。

\textrm{VASP}首先完成电子结构的最小化，得到超晶胞中的基态。然后\textrm{VASP}分别沿着\textrm{Cartesian}坐标的三个方向每次偏移一个原子，再计算对应的\textrm{Kohn-Sham}相互作用势的改变。计算结果写入文件\textsc{phelel\_params.hd5}。%注意，当前计算中只用$\Gamma$点，所以计算中采用\textrm{vasp\_gam}版本完成计算。

因为设置了计算控制参数\texttt{ELPH\_POT\_GENERATE}，所以会产生输出文件\textrm{CONTAR\_ELPH}。该文件类似于\textrm{CONTCAR}文件，不过文件内容是初基原胞的晶体结构信息。对照\textrm{POSCAR}文件和\textrm{CONTCAR\_ELPH}文件的内容。
fbox{\mathrm{vim -O POSCAR CONTCAR\_ELPH}}

后续的电子-声子计算中，\textrm{CONTCAR\_ELPH}就是作为输入文件\textrm{POSCAR}来使用。如果用户不希望\textrm{CONTCAR\_ELPH}保存的是初基原胞的结构，只要在\textrm{INCAR}文件中设置计算控制参数\texttt{ELPH\_POT\_LATTICE}。

由超晶胞体系的\textrm{VASP}计算得到的电子-声子相互作用势戏被写入文件\textsc{phelele\_params.hd5}，该相互作用势是基于一个$\vec k$空间网格，对元胞计算得到的。根据计算控制参数\texttt{ELPH\_POT\_FFT\_MESH}的设定，电子-声子相互作用势是在精密\textrm{FFT}网格(超晶胞)上计算的，但输出时作为\textrm{FFT}网格(元胞)形式完成的。默认条件下，在电子-声子计算过程中，计算控制参数\texttt{ELPH\_POT\_FFT\_MESH}网格必须是和\textrm{FFT}网格(元胞)有完全相同空间划分。

对于元胞计算，\textrm{VASP}选定的\textrm{FFT}网格点的影响因素有
\begin{itemize}
	\item 首先，\textrm{FFT}网格点依赖电子截断能量\textrmENCUT
	\item 通过计算控制参数\texttt{PREC}确定\textrm{FFT}网格的数目
	\item 直接设置计算控制参数\texttt{NGX},\texttt{NGY}和\texttt{NGZ}
\end{itemize}
所以如果计算过程中，对超晶胞和元胞计算，没有改变计算控制参数\texttt{ENCUT}和\texttt{PREC}，那么这两套网格点将会自动匹配，不需要进一步的输入处理。

不过，在有些情况下，由精密\textrm{FFT}网格(对超晶胞)计算得到电子-声子相互作用势，再通过默认的\textrm{FFT}网格(对元胞)表示，则并不一定完全收敛，包括其他的电子-声子性质(如带隙重整化)也有类似情形。如果试图通过改变\texttt{ENCUT}来定义一套新的\textrm{FFT}网格(对元胞)过于复杂，留待后续讨论。相反，通过计算控制参数\texttt{ELPH\_POT\_FFT\_MESH}，可以在新的一轮超晶胞计算过程中再定义的一套新\textrm{FFT}网格，显然，代价是需要重新完成一次完整的超晶胞计算。

